\scanpage{354}
% top centre
\mspage{270}
% lines 1–5
\inserted{(resp. par $\varinjlim$)} de cardinal $\leq\pi$, donc par 4.11.1 $M$ est stable par ces $\unsure{\overline{\mathrm m}^{\mathrm{es}}|\overline{\mathrm m}\mathrm s}$ \inserted{types de} limites inductives. \deleted{Il s'ensuit} D'autre part il est stable par $\varinjlim$ \unsure{grdes|gdes} devant $\pi$. Il suffit donc de prouver le

% lines 6–11
\uline{Lemme} 4.\unsure{14|11}.3. Soit $M$ une catégorie, $\pi$ un cardinal. Supposons que $M$ soit stable par $\varinjlim$ filtrantes de cardinal $\leq\pi$ \inserted{(resp par \deleted{\unsure{$\varinjlim$ de}} $\varinjlim$ de cardinal $\leq\pi$, resp \unsure{par sommes|sommes} de cardinal $\leq\pi$)}, \uuline{et} par $\varinjlim$ \unsure{grdes|gdes} devant $\pi$. Alors $M$ est stable par \inserted{petites} $\varinjlim$ filtrantes, (resp. par \inserted{petites $\varinjlim$, resp. par} \deleted{petites} petites sommes).

% line 12
Ce \inserted{lemme se} prouve comme 4.4.14 page 97.

% lines 13–14
Reste la réciproque dans 4.\unsure{11|10}.2, qui résulte de l'assertion similaire dans 4.11.1.

% lines 15–19
\uline{Corollaire} 4.11.4. Quand $M$\inserted{$=\operatorname{Ind}_\pi C$} est-il stable par sommes ? Si $\exists$ ordinal $\alpha$ tel que $\mathrm{card}\,\alpha\leq\pi$ et que $C$ soit stable par $\varinjlim$ de type $I_\alpha$, alors il suffit que $C$ soit stable par sommes de cardinal $\leq\pi$. Cette condition est aussi nécessaire \unsure{{}|si} $C$ est karoubienne.
% lower left beside last corollary
% slanted
\margin{Pour \unsure{autres} \unsure{variantes|variations} de 4.11.1, cf 4.13.19 (p. 340) et corollaires.}

% bottom line
Se prouve comme 4.11.2, en utilisant 4.11.3.
% archive sheet number: 353
